\section{辅政、外戚与新朝的禅代}
\label{sec:western-han-regents-wang-mang}

武帝死后，西汉并未立刻回到一个由成年皇帝直接统御的秩序。前八七年，年少的刘弗陵即位，霍光等受遗诏辅政；辅政安排能够延续诏令、禁军和人事，却也把“谁代皇帝决定”的问题留在中枢。前八一年围绕盐铁、专营、边防与财政的廷议，保存于《盐铁论》及《汉书》的不同层次材料中。它不能被简化为朝廷已在一次辩论后废止或保留了所有政策，却显示武帝时代留下的军事和财政负担，已成为官员必须公开争论的政务。

\subsection{辅政不是无声的过渡}

前七四年昭帝无嗣而死，霍光先迎立昌邑王刘贺，旋即废之，改立刘病已为宣帝。废立的日期和结果可靠，昌邑王被归纳为“失德”的种种罪状则主要来自证明废立正当性的叙事，不能逐条改写成现场记录。此事更直观地说明，辅政大臣已拥有选择皇帝的实际能力。皇统没有因此中断，却必须通过功臣、外戚、宫廷官员和宗室都能接受的安排，重新取得承认。

霍光前六八年去世后，宣帝逐步亲政并削弱霍氏。宫廷婚姻、许皇后的死亡以及霍氏遭清算的先后与责任仍有争议，不宜把政治变化缩减为一桩后宫恩怨。较稳的线索是，受遗诏辅政所积累的权力并不能自动回到皇帝手中；它需要通过人事、司法和宗室关系被重新分配。宣帝朝对狱讼与吏治的诏令，和赵充国在羌地所主张的屯田、招抚、分化并行，也提示“守成”仍要依靠具体的官员、地方道路和边郡资源，而非仅靠少征或不征。

前六〇年郑吉受命为西域都护，前五一年呼韩邪单于入朝。都护的治所、管辖范围和南匈奴的依附程度都有需要逐案辨析的地方；不能因为有了一个官名或一次朝见，便把绿洲、草原与边郡画成完全受控的内地。它们至少表明，宣帝朝以使节、赏赐、屯戍和地方军事相结合，维持了一种不同于武帝大规模远征的边疆关系。

\subsection{王氏进入中枢}

前四九年宣帝去世，元帝继位；至成帝朝，王凤长期领尚书事，王氏外戚逐渐形成中枢集团。后宫、外戚与宗室的竞争并非成帝一朝突然产生：帝嗣不定、官爵分配和皇太后家族的可用人脉，长期为它提供了条件。前一六年王莽受封新都侯，前八年出任大司马。关于他的孝行形象、谦退姿态和个人谋划，主要来自后来高度道德化的传记，不能把这种人物塑像当作权力上升的全部解释；可以确知的是，他已进入能够影响官僚任用与朝廷议程的位置。

前一年哀帝无嗣而死，王政君召王莽入朝，中山王刘衎继位为平帝。幼帝与外戚辅政的组合，使皇权名义和实际决策再度分开。元始二年所存户口统计，是一份来自帝国行政申报的珍贵尺度，却不能直接转换为社会总人口或税役负担的透明数字；它所显示的是，朝廷仍努力以郡国簿籍来表达对广阔人口与土地的把握。

\subsection{从居摄到新}

平帝于五年去世，王莽立孺子婴为皇太子并居摄；六年以后，他以诏令调整官名、处置反对者，八年又接受符命和劝进，九年受禅称帝，改国号新。居摄、符命和禅代均可在《汉书》〈王莽传〉、帝纪与官表中找到年序，平帝死因、符命如何制造、群臣是否出于自愿则不能由这些文本直接裁定。新莽钱币和官印材料能为名号、制度变更提供物质锚点，也不能说明每一道新令在地方实行到什么程度。

王莽的受禅不是西汉制度失效后的空白，而是西汉长期解决过、又不断留下的问题被重新聚合：辅政者可以支配人事，外戚可以组织朝廷网络，幼主的继承可以被礼仪包装，而帝国的簿籍与官名仍能被用来宣称新的秩序。新朝能否在郡县、边地和家庭中真正取得服从，是另一组尚待逐地检验的问题；至少在九年，汉的皇帝继承线已被一场有制度外衣的权力转移所截断。

\subsection{史料的取舍}

《汉书》对霍光、王氏与王莽提供了最连续的主干，同时也最强烈地以东汉的正统立场组织褒贬。《汉纪》可帮助核对年序，却是后出的编年叙述；《盐铁论》保留了政策语言和论辩形式，文本的编纂层次使它不能直接替代前八一年的朝廷实录。西域传、匈奴传、钱币、官印和户籍材料分别照出边疆、制度和行政的一角。它们共同要求叙事把废立、外戚、边防和财政放在同一时代的资源约束中理解，而不将任何一位辅政者或皇太后的道德形象当作唯一原因。
